% Ad Atticum 10.10 --- headnote
% ad-atticum-10-10, 3 May 49 BC, Cumae

\textit{Cicero to Atticus, written from the Cuman villa
on the same day as Att.\ 10.9 --- the fifth day before
the Nones of May 49 BC, 3 May (the manuscript dateline:
\textit{Scr.\ in Cumano v Non.\ Mai.\ a.\ 705 (49)}). The
letter is built around the text of a reply Cicero has
just received from Mark Antony, who has been put in
charge of the Italian coast and given Caesar's instruction
to let no one out. Cicero had written repeatedly,
disclaiming any move against Caesar's interests, naming
his son-in-law Dolabella, naming the old friendship,
saying he could have gone with Pompey if he wished and
that for himself he simply wanted to be away from the
lictors he was being made to attend. The Latin opens
with a flash of bitter self-rebuke (\textit{me caecum
qui haec ante non viderim}): Cicero now sees that he
has been answered with deliberate, smiling refusal.
The Greek tag \textit{[Greek: parainetik\=os]} ---
``in an exhortative manner'' --- is dropped in as
hostile shorthand: that is, with the tone of a friend
giving advice, but a refusal all the same. Section 2 is
the quoted letter itself, the only stretch in Cicero's
correspondence in which Mark Antony speaks directly:
courteous, irreproachable, and entirely closing the
door. The middle is its sting: he has no authority to
let Cicero go, only the Caesarian remit to keep
everyone in Italy; if Cicero wants to leave, he must
write to Caesar.}

\textit{Section 3 is the response in code: \textit{habes
\textit{[Greek: skytal\=en Lak\=onik\=en]}} --- ``there
is your Spartan scytale,'' the cipher staff. Antony's
courtesy reads, in cipher, as a refusal; Cicero will
take the meeting (Antony is due that very evening),
will make a show of unhurried compliance, will say he
will write to Caesar --- and meanwhile, in secret, with
the fewest possible companions, will slip out. The
allusive Greek line \textit{[Greek: sunes ho toi
leg\=o]} (``mark what I am telling you'') is a
Pindaric tag dropped in to flag what is being said
beneath the surface. Two daggered cruxes in the
manuscript (\dag{}audeam\dag{}, \dag{}carti\dag{}) sit
inside this same sentence; the run of the thought is
intact even where the text is corrupt. Section 4 turns
to Atticus's health (the strangury, treat it while it
is still at its \textit{[Greek: arch\=e]}) and to news
out of Massilia. Section 5 is the unforgettable image
of Antony on the move with the actress Cytheris in an
open litter like a wife, seven more linked litters of
mistresses and boyfriends behind: ``See what a shameful
death we are perishing by'' --- and the resolve, even
in a skiff if no ship can be had, to tear free. Section
6 returns one last time to the long anxiety about his
nephew Quintus: a remarkable wit but a still-unformed
character \textit{[Greek: an\=ethopoi\=eton]}, turned
away from his own people; the character must be tended
to.}
